\section{Review Checklist}

\subsection{Before Writing}

Before replacing this template with a real topic, define the paper's core claim in one sentence. A claim such as ``this implementation reduces median request latency under a fixed workload'' is easier to evaluate than a broad theme such as ``performance matters.'' The claim determines which background material is necessary, which experiments should be included, and which results are relevant.

The writer should also identify the expected reader. A classmate may need more implementation context, while a reviewer or instructor may expect clearer assumptions and stronger references. The template does not force one level of detail, but it provides enough structure to support either emphasis.

\subsection{Before Submission}

Use the following checks before considering the paper complete:
\begin{enumerate}
\item The abstract states the problem, method, and result.
\item Every nontrivial technical claim has a citation, derivation, or experiment.
\item All figures and tables are referenced from the text before or near their placement.
\item All citations resolve and the bibliography appears in IEEE style.
\item The repository link points to the final public or private repository expected by the assignment.
\item The statement section reflects the required academic integrity wording.
\item The PDF was regenerated from a clean local compile, not from a stale build artifact.
\end{enumerate}

This checklist is intentionally practical. Formatting errors are usually easy to fix, but missing claims, unclear assumptions, and unsupported results are harder to repair near a deadline. A short review pass focused on those issues improves the paper more than repeated cosmetic editing.

\subsection{Version Control Hygiene}

Generated LaTeX files should not be committed unless the assignment, venue, or project policy explicitly asks for the final PDF in the repository. The source of truth is the combination of \texttt{main.tex}, section files, references, figures, and README instructions. Keeping generated files out of commits makes it easier to review the actual writing and prevents large binary diffs from hiding meaningful changes.

When working over multiple sessions, compile before stopping and check the git status. A clean build confirms that the next editing session starts from a known-good source tree. If a section is incomplete, leave a clear placeholder paragraph rather than a broken command or missing file input.
